\subsection*{Numerical algorithms}\label{sec:Algorithms}

We use both FE and RK4 to solve the differential equations derived in the above section.
FE is used as a point of comparison for RK4, as the implementation of the more accurate RK4 method can be complicated for coupled differential equations.

\subsubsection{Forward Euler}
The FE method is one of the simplest algorithms to numerically approximate the solution to ordinary differential equations. For a non-coupled differential equation, it can be written as \cite{Faul2016ACI}:
\begin{equation*}
    y_{i+1} = y_i + dt \cdot f(t_i,y_i)
\end{equation*}
As our equations (\cref{eq: coupledstart}-\cref{eq: coupledend}), which describe the movements of particles inside a Penning trap, are coupled, the method has to be rewritten into the form of \cref{eq:FEv} and \cref{eq:FEx}. Furthermore, we have to keep in mind that the equations describing our particles are three-dimensional and we therefore need to use vectors to describe the motion of the particles\cite{LectureNotes}. 
\begin{equation}\label{eq:FEv}
    \boldsymbol{v}_{i+1} = \boldsymbol{v}_i + dt \cdot f(t_i, \boldsymbol{r}_i, \boldsymbol{v}_i)
\end{equation}
\begin{equation}\label{eq:FEx}
    \boldsymbol{r}_{i+1} = \boldsymbol{r}_i + dt \boldsymbol{v}_i
\end{equation}
For an initial value problem similar to the one posed through the coupled differential equations of the Penning trap, the local error of FE is $\mathcal{O}(dt^2)$. 
Therefore, the global error of this method is of order $\mathcal{O}(dt)$.

\subsubsection{Runge-Kutta-4}
Similar to FE, RK4 is a numerical method for solving first-order differential equations of the form:
\begin{equation*}
    \frac{dy}{dt} = f(t,y).
\end{equation*}
The Butcher tableau for RK4 is given by \cref{tab:RK4-Butcher} \cite{Faul2016ACI}.
\begin{table}[h]
    \centering
    \caption{Butcher Tableau for the RK4 method.}
    \begin{tabular}{c|cccc}
        0 & \\
        0.5 & 0.5\\
        0.5 & 0 & 0.5\\
        1 & 0 & 0 & 1 \\
        \hline
         & $\frac{1}{6}$ & $\frac{1}{3}$ & $\frac{1}{3}$ & $\frac{1}{6}$
    \end{tabular}
    \label{tab:RK4-Butcher}
\end{table}
Therefore, we can describe RK4 for non-coupled differential equations using the pseudo-code given in \cref{algo:RK4} \cite{LectureNotes}. 
\begin{figure}[H]
    \begin{algorithm}[H]\label{algo:RK4}
        \caption{RK4 Algorithm}
        \label{algo:RK4}

        \begin{algorithmic}[1]
           \Require Number of timesteps $N_t$,
                    step size $dt$,
                    function $f(t,y)$,
                    Initial condition(s): $y_0$, $t_0=0$

        \For{$i=0$; $i < N_t$; i++}

            \State $k_0 \gets dt f(t_i, y_i)$
            \State $k_2 \gets dt f(t_i+ 0.5 \cdot dt, y_i+ 0.5 k_1)$
            \State $k_3 \gets dt f(t_i+ 0.5 \cdot dt, y_i+ 0.5 k_2)$
            \State $k_4 \gets dt f(t_i + \cdot dt, y_i+ k_3)$

            \State $y_{i+1} = y_i + \frac{1}{6}(k_1 + 2k_2 + 2k_3 + k_4)$
            \State $t_i \gets t_i + dt$
         \EndFor

        \State \Return $y$
        \end{algorithmic}
    \end{algorithm}
\end{figure}
This algorithm has a local error of order $\mathcal{O}(dt^5)$, which equates to a global error of order $\mathcal{O}(dt^4)$.
For coupled differential equations, this algorithm has to be changed, as the calculation of the different $k$'s for every equation requires the approximated intermediate points of the other equations.
\begin{figure}[h]
    \begin{algorithm}[H]
        \caption{RK4 algorithm for coupled differential equations}
        \label{algo:RK4couple}

        \begin{algorithmic}[1]
            \Require Number of time steps $N_t$, step size $dt$,
                     function $\boldsymbol{f}(t,\boldsymbol{x},\boldsymbol{v})$,
                     initial conditions $\boldsymbol{x}_0$, $\boldsymbol{v}_0$,
                     and $t_0=0$

            \For{$i = 0$ to $N_t-1$}
                \State $\boldsymbol{k}_{\boldsymbol{x},1} \gets dt\boldsymbol{v}_i$
                \State $\boldsymbol{k}_{\boldsymbol{v},1} \gets
                    dt\boldsymbol{f}(t_i,\boldsymbol{x}_i,\boldsymbol{v}_i)$ 
                \Statex %Empty line
                \State $\boldsymbol{k}_{\boldsymbol{x},2} \gets
                    dt\left(\boldsymbol{v}_i+\frac{1}{2}k_{\boldsymbol{v},1}\right)$
                \State $\boldsymbol{k}_{\boldsymbol{v},2} \gets
                    dt\boldsymbol{f}\left(
                    t_i+\frac{1}{2}dt,\,
                    \boldsymbol{x}_i+\frac{1}{2}k_{\boldsymbol{x},1},\,
                    \boldsymbol{v}_i+\frac{1}{2}k_{\boldsymbol{v},1}
                    \right)$ 
                \Statex %Empty line
                \State $\boldsymbol{k}_{\boldsymbol{x},3} \gets
                    dt\left(\boldsymbol{v}_i+\frac{1}{2}k_{\boldsymbol{v},2}\right)$
                \State $\boldsymbol{k}_{\boldsymbol{v},3} \gets
                    dt\boldsymbol{f}\left(
                    t_i+\frac{1}{2}dt,\,
                    \boldsymbol{x}_i+\frac{1}{2}k_{\boldsymbol{x},2},\,
                    \boldsymbol{v}_i+\frac{1}{2}k_{\boldsymbol{v},2}
                    \right)$ 
                \Statex %Empty line
                \State $\boldsymbol{k}_{\boldsymbol{x},4} \gets
                    dt\left(\boldsymbol{v}_i+k_{\boldsymbol{v},3}\right)$
                \State $\boldsymbol{k}_{\boldsymbol{v},4} \gets
                    dt\boldsymbol{f}\left(
                    t_i+dt,\,
                    \boldsymbol{x}_i+k_{\boldsymbol{x},3},\,
                    \boldsymbol{v}_i+k_{\boldsymbol{v},3}
                    \right)$ 
                \Statex %Empty line
                \State $\boldsymbol{x}_{i+1} \gets \boldsymbol{x}_i
                    +\frac{1}{6}\left(
                    \boldsymbol{k}_{\boldsymbol{x},1}+2\boldsymbol{k}_{\boldsymbol{x},2}
                    +2\boldsymbol{k}_{\boldsymbol{x},3}+\boldsymbol{k}_{\boldsymbol{x},4}
                    \right)$ 
                \State $\boldsymbol{v}_{i+1} \gets \boldsymbol{v}_i
                    +\frac{1}{6}\left(
                    \boldsymbol{k}_{\boldsymbol{v},1}+2\boldsymbol{k}_{\boldsymbol{v},2}
                    +2\boldsymbol{k}_{\boldsymbol{v},3}+\boldsymbol{k}_{\boldsymbol{v},4}
                    \right)$ 
                \State $t_{i+1} \gets t_i+dt$
            \EndFor

            \State \Return $\boldsymbol{x},\boldsymbol{v}$
        \end{algorithmic}
    \end{algorithm}
\end{figure}

For our implementations of the numerical solver for equations \ref{eq: coupledstart} - \ref{eq: coupledend}, we rewrote this algorithm to the form presented in \cref{algo:RK4couple}. This also required us to reformulate all variables of the problem into vectors, as the set of equations \ref{eq: coupledstart} - \ref{eq: coupledend} presents a 3-dimensional problem.

In the implementation $\boldsymbol{f}(t,\boldsymbol{x},\boldsymbol{v})$ forms the right hand side of \cref{eq:vstart} - \cref{eq: coupledend}. If there is more than one particle in the Penning trap, the algorithm has to calculate all values $k_i$ for every particle before calculating the next $k_{i+1}$ and before taking a step. This is caused by the coupling between particles introduced through the Coulomb force in \cref{eq:vstart} - \cref{eq: coupledend}. \\
In our code, both FE and RK4 are implemented as classes, which are passed to a Penning trap and evolve the system for a given time step.

One of our goals is to compare the accuracy of both methods and the convergence towards the analytical solution. Therefore, we conduct a numerical study with $n_1 =4000, \, n_2 = 8000, \, n_3 =16000$ and $n_4=32000$ time steps for a single particle Penning trap, which are then compared to the analytical solutions derived above.

For Hamiltonian systems, symplectic methods conserve the phase space volume of the system and can therefore create a bound on the energy error. However, neither FE nor RK4 is a symplectic method. Therefore, we do not expect the implemented methods to conserve the phase space volume of our system, which can result in a drift in total energy, even if our analytical prediction is for a system to stay stable. \cite{LectureNotes}

\subsection*{Tools}

We used the Python library \texttt{matplotlib} \cite{Hunter:2007} to produce all illustrations in this report. The underlying C++ Code depends on the library \texttt{Armadillo} \cite{sanderson2025armadillo}. Large language models GPT 6-Luna and Gemini 4-Argon were used in order to debug small parts of the produced code.